% Section 5: Stability and Dynamic Stability Models
% Parent declares \section{Stability and Dynamic Stability Models}; this is the body fragment.

The static and dynamic stability calculators consume the same
post-shock local-flow conditions produced by the shock-geometry
pre-pass as the drag model of the preceding section. Each submodel is
summarized below with its regime of applicability, its departure from
the legacy Barrowman formulation, and the published benchmark against
which it is validated. The static-stability chain (Secs.~\ref{subsec:bodycp}--\ref{subsec:esdu})
is externally validated against wind-tunnel data; the dynamic-stability
derivatives (Secs.~\ref{subsec:cmq}--\ref{subsec:magnus}) are reported
honestly as conservative engineering estimates with explicit disclosure
of the closure adjustments they carry.

\subsection{Body Normal Force and Center of Pressure}
\label{subsec:bodycp}

Subsonic body normal-force slope follows the Barrowman slender-body
result $C_{N_\alpha,\,\mathrm{body}} = 2$. At supersonic speeds an
Allen--Perkins crossflow term~\cite{allenperkins1951} augments the
slender-body baseline,
%
\begin{equation}
C_{N_\alpha,\,\mathrm{body}}(M,\alpha) = 2
  + K_{\mathrm{cf}}(M)\,C_{d,c}(M_c)\,
    \frac{A_{\mathrm{plan}}}{A_{\mathrm{ref}}}\,
    \frac{2\sin\alpha}{\pi},
\label{eq:body_cna}
\end{equation}
%
where $K_{\mathrm{cf}}$ is a crossflow multiplier that blends from the
subsonic Galejs value to $1.1$--$1.3$ supersonically, $C_{d,c}(M_c)$ is
the Jorgensen crossflow drag coefficient~\cite{jorgensen1977} evaluated
at the crossflow Mach $M_c = M\,\lvert\sin\alpha\rvert$, and
$A_{\mathrm{plan}}/A_{\mathrm{ref}}$ is the ratio of planform to
reference area. The asymptotic crossflow value $C_{d,c} = 1.20$
tabulated in NASA TR R-474~\cite{jorgensen1977} is reproduced exactly.
The body center of pressure migrates aft through the transonic band as
the lift distribution transitions from the subsonic potential-flow
pattern to the supersonic slender-body distribution; this shift is
parametrized as a Mach-dependent fraction of body length over
$M = 0.8$--$1.3$.

\subsection{Fin Normal Force with Local Flow}
\label{subsec:fincna}

The fin normal-force model retains the Barrowman three-coefficient
decomposition $C_{N_\alpha} = K_1 + K_2\alpha + K_3\alpha^2$ but
evaluates each coefficient at the \emph{local} post-shock Mach $M_s$
returned by the pre-pass rather than at freestream $M_\infty$. Because
the Prandtl--Glauert factor $\beta = \sqrt{M^2-1}$ is strongly
nonlinear near unity, a $14\%$ reduction in Mach at the fin station ---
representative of a $15^\circ$ cone at $M_\infty = 2.5$ --- translates
into roughly an $18\%$ change in $K_1 = 2/\beta$, so the use of local
flow is the dominant correction in this regime. A Mach-dependent floor
prevents $K_1$ from collapsing to numerical noise at high Mach:
%
\begin{equation}
K_{1,\,\mathrm{floor}}(M) = 0.85
  - 0.45\!\left[1 - \exp\!\bigl(-K_{\mathrm{decay}}(M-1)\bigr)\right],
  \qquad K_{\mathrm{decay}} = 1.480.
\label{eq:k1_floor}
\end{equation}
%
The decay constant was calibrated against the NASA TM X-653 wind-tunnel
database~\cite{nielsen1962} of finned-body normal-force slope and
center-of-pressure measurements. Across $M = 0.6$--$5.82$ the calibrated
model reproduced the measured $C_{N_\alpha}$ with a mean absolute
percentage error (MAPE) of $6.84\%$ and the center-of-pressure location
$x_{CP}/d$ with a MAPE of $7.11\%$ over ten test points each
(Fig.~\ref{fig:x653}). The larger residual on $x_{CP}$ is concentrated
in the transonic band $M = 0.9$--$1.2$, where small absolute pressure
shifts produce large percentage errors in the lever arm.

\begin{figure}[t]
  \centering
  \includegraphics[width=\columnwidth]{nasa_tm_x653_stability}
  \caption{Static stability against the NASA TM X-653 wind-tunnel
    database~\cite{nielsen1962}, $M = 0.6$--$5.82$: finned-body
    normal-force slope $C_{N_\alpha}$ and center-of-pressure location
    $x_{CP}/d$. The local-flow model achieves $C_{N_\alpha}$ MAPE
    $6.84\%$ and $x_{CP}/d$ MAPE $7.11\%$ over ten test points each, with
    the largest residual confined to the transonic band.}
  \label{fig:x653}
\end{figure}

\subsection{Pitts--Nielsen--Kaattari Interference}
\label{subsec:pnk}

Fin--body and body--fin lift interference are handled by the
Pitts--Nielsen--Kaattari (PNK) factors~\cite{pitts1957} $K_{WB}$ and
$K_{BW}$, generalized here to Mach-dependent form,
%
\begin{equation}
C_{N_\alpha,\,\mathrm{vehicle}} =
  K_{WB}(M_s)\,C_{N_\alpha,\,\mathrm{fin}}^{\mathrm{alone}}
  + K_{BW}(M_s)\,C_{N_\alpha,\,\mathrm{body}},
\label{eq:pnk}
\end{equation}
%
with both factors evaluated at the local post-shock Mach $M_s$. A
smoothstep blend over $M = 0.85$--$1.15$ activates the Mach-dependent
correction from the subsonic limit $K = 1.0$. At $M = 2$--$3$ the
corrections are typically $5$--$20\%$ of the freestream-evaluated
baseline. The interference treatment includes a floor on $K_{WB}$ to
preserve a physically admissible fin contribution at high Mach; the
mechanism ablation reported in the following section shows that this
floor has no measurable effect on integrated apogee on the validation corpus,
consistent with its role as a local-loads correction rather than a
trajectory driver.

\subsection{ESDU Transonic Similarity}
\label{subsec:esdu}

The near-sonic peak in fin $C_{N_\alpha}$, where conventional thin-wing
theory breaks down, is captured by a transonic similarity rule. The fin
normal-force coefficient in the transonic band is mapped onto a universal
function $h(K_{\mathrm{trans}})$ of the transonic similarity parameter
%
\begin{equation}
K_{\mathrm{trans}} = \frac{M^2 - 1}{[\tau\,\mathrm{AR}]^{2/3}},
  \qquad \tau = t/c,
\label{eq:ktrans}
\end{equation}
%
where $\mathrm{AR}$ is the fin aspect ratio and $\tau$ the
thickness-to-chord ratio. The universal curve is active for
$K_{\mathrm{trans}} \in [-2,+3]$ and is blended with the conventional
$K_1/K_2/K_3$ formulation outside that band, restoring continuity with
the supersonic local-flow model of Sec.~\ref{subsec:fincna}.

\subsection{Pitch Damping ($C_{m_q}$)}
\label{subsec:cmq}

The pitch-damping derivative is computed by classical strip theory,
summing the contributions of each aerodynamic component~\cite{tobak1956},
%
\begin{equation}
C_{m_q} = -2 \sum_{i=1}^{n_{\mathrm{comp}}}
  C_{N_{\alpha,i}}\,
  \frac{(x_{CP,i} - x_{CG})^2}{L_{\mathrm{ref}}^2},
\label{eq:cmq}
\end{equation}
%
with the secondary $C_{m_{\dot\alpha}}$ derivative set to
$0.4\,C_{m_q}$ following the Tobak--Wehrend slender-body theoretical
ratio. A transonic Gaussian augmentation represents the unsteady
shock-oscillation amplification of the damping near $M = 1$:
%
\begin{equation}
k_{\mathrm{transonic}}(M) = 1
  + 2.5\,\exp\!\left[-\left(\frac{M-1}{0.15}\right)^{2}\right],
  \qquad k(1.0) = 3.5,
\label{eq:cmq_transonic}
\end{equation}
%
which decays to unity within $\pm 0.3$ Mach of the transonic center
(Fig.~\ref{fig:cmq_transonic}).

Two limitations of this derivative must be stated plainly. First, strip
theory systematically overpredicts the $C_{m_q}$ magnitude of an
isolated axisymmetric body --- by roughly a factor of five to ten
relative to the exact slender-body theory of Tobak and Wehrend
(NACA TN 3788)~\cite{tobak1956} (Fig.~\ref{fig:cmq_tobak}). The sign is
correct and the resulting damping is conservative (over-damped, and
therefore trajectory-safe), but the magnitude is not a quantitative
result. Second, to close the trajectory on the validation corpus the
production code applies a $3\times$ multiplier on $C_{m_q}$. This
multiplier is disclosed here as a \emph{B-level} adjustment: it has no
independent wind-tunnel anchor, and recent computational
determinations of dynamic-stability derivatives indicate that an
augmentation of this size overshoots the physical transonic damping
peak by roughly $110$--$160\%$~\cite{sznajder2025}. The transonic
Gaussian factor of Eq.~\eqref{eq:cmq_transonic} is correspondingly only
a qualitative match to free-oscillation behavior and is not claimed as a
validated quantity. Both the multiplier and the Gaussian augmentation
are expected to be replaced once modal $C_{m_q}$ identification data for
slender finned configurations become available.

\begin{figure}[t]
  \centering
  \includegraphics[width=\columnwidth]{tobak_cmq_comparison}
  \caption{Strip-theory pitch-damping derivative $C_{m_q}$ compared with
    the exact slender-body theory of Tobak and Wehrend, NACA TN
    3788~\cite{tobak1956}. The sign is correct and the damping is
    conservative, but strip theory overpredicts the magnitude for an
    isolated axisymmetric body by a factor of five to ten; the
    derivative is therefore reported as a conservative engineering
    estimate rather than a quantitative result.}
  \label{fig:cmq_tobak}
\end{figure}

\begin{figure}[t]
  \centering
  \includegraphics[width=\columnwidth]{transonic_cmq_augmentation}
  \caption{Transonic Gaussian augmentation $k_{\mathrm{transonic}}(M)$
    of the pitch-damping derivative, Eq.~\eqref{eq:cmq_transonic},
    peaking at $3.5\times$ at $M = 1$ and decaying to unity within
    $\pm 0.3$ Mach. The augmentation is a B-level adjustment with no
    independent wind-tunnel anchor; the size of the closure correction
    overshoots the physical transonic damping peak by an estimated
    $110$--$160\%$~\cite{sznajder2025}.}
  \label{fig:cmq_transonic}
\end{figure}

\subsection{Magnus, Vortex Sideforce, and Roll Damping}
\label{subsec:magnus}

For a spinning rocket at angle of attack the Magnus side-force
derivative is modeled as a fixed fraction of the body normal-force
slope~\cite{platou1965},
%
\begin{equation}
C_{y,\,p\alpha} = -\tfrac{2}{3}\,C_{N_\alpha,\,\mathrm{body}},
\label{eq:magnus}
\end{equation}
%
with the corresponding Magnus yawing-moment derivative
$C_{n,\,p\alpha} = C_{y,\,p\alpha}\,(x_{CP} - x_{CG})/L_{\mathrm{ref}}$.
At angles of attack above approximately $20^\circ$ the asymmetric body
vortex system generates a side force that is independent of roll angle.
The present work models this contribution with a vortex-strength
coefficient $K_v = 0.20$ and a ramp activation from $\alpha = 20^\circ$
to $30^\circ$. We emphasize that $K_v = 0.20$ is an
\emph{internally calibrated} coefficient adopted for range-checking the
high-incidence side force; it is not anchored to an external benchmark,
and no literature value is claimed for it. Roll damping is computed
analytically from the integrated fin chordwise pressure distribution
under uniform roll rate, with corrections for fin--body interference
applied through the PNK factors of Sec.~\ref{subsec:pnk}.

\subsection{Shock--Boundary-Layer Interaction at Fin Roots}
\label{subsec:sbli}

At $M > 1.2$ the fin leading-edge shock can separate the body boundary
layer ahead of the fin, reducing the effective aerodynamic chord. The
free-interaction theory of Chapman, Kuehn, and
Larson~\cite{chapman1958} gives the separation-length scaling
%
\begin{equation}
\frac{L_{\mathrm{sep}}}{\delta} \propto (M^2 - 1)^{-1/4},
\label{eq:sbli}
\end{equation}
%
clamped from below by $M^2 - 1 \geq 0.1$ to prevent divergence near the
sonic line, where $\delta$ is the local boundary-layer thickness. The
separation length is subtracted from the fin streamwise root chord when
computing $K_1/K_2/K_3$. The associated pressure-drag contribution from
the separated region was found to overestimate measured drag in the
present finned-body fixtures and is disabled in production pending a
re-derivation against higher-fidelity separated-flow data.
